\frfilename{mt255.tex}
\versiondate{3.7.08}
\copyrightdate{2000}

\def\chaptername{Ukuran produk}
\def\sectionname{Konvolusi fungsi}

\newsection{255}

Saya mengkhususkan satu bagian untuk suatu konstruksi yang sangat penting -- dan
khususnya akan sangat berguna dalam Bab 27 dan 28 -- dan juga dapat
dipandang sebagai serangkaian latihan atas pembahasan sejauh ini.

Saya merasa sulit menentukan seberapa banyak pengulangan yang patut diberikan dalam
bagian ini, karena ungkapan terpadu yang alami bagi gagasan-gagasan ini terdapat dalam
teori grup
topologis, dan saya rasa kita belum siap untuk teori umum itu (saya
akan membahasnya dalam Bab 44 pada Jilid 4).   Grup-grup yang kita perlukan untuk jilid ini adalah

\qquad $\Bbb R$;

\qquad $\BbbR^r$, untuk $r\ge 2$;

\qquad $S^1=\{z:z\in\Bbb C,\,|z|=1\}$, `grup lingkaran';

\qquad $\Bbb Z$, grup bilangan bulat.

\noindent Semua gagasan tersebut sudah muncul dalam teori konvolusi pada
$\Bbb R$, dan karena itu saya akan menyajikan bahan ini dalam bentuk yang relatif
terperinci, sebelum menguraikan secara ringkas bentuk-bentuk yang sesuai untuk grup
$\BbbR^r$ dan $S^1$ (atau $\ocint{-\pi,\pi}$);  menurut saya, $\Bbb Z$ dapat
dengan aman diserahkan kepada latihan.

\leader{255A}{}\cmmnt{ Karena ini buku tentang teori ukuran, kiranya
tepat jika saya menekankan, sebagai dasar teori
konvolusi, beberapa isomorfisme ruang ukur.

\medskip

\noindent}{\bf Teorema} Misalkan $\mu$ ukuran Lebesgue pada $\Bbb R$ dan
$\mu_2$ ukuran Lebesgue pada $\BbbR^2$;  tuliskan $\Sigma$, $\Sigma_2$ untuk
domain masing-masing.

(a) Untuk sebarang $a\in\Bbb R$, pemetaan $x\mapsto a+x:\Bbb R\to\Bbb R$ merupakan
automorfisme ruang ukur dari $(\Bbb R,\Sigma,\mu)$.

(b) Pemetaan $x\mapsto -x:\Bbb R\to\Bbb R$ merupakan automorfisme ruang
ukur dari $(\Bbb R,\Sigma,\mu)$.

(c) Untuk sebarang $a\in\Bbb R$, pemetaan $x\mapsto a-x:\Bbb R\to\Bbb R$ merupakan
automorfisme ruang ukur dari $(\Bbb R,\Sigma,\mu)$.

(d) Pemetaan $(x,y)\mapsto(x+y,y):\BbbR^2\to\BbbR^2$ merupakan automorfisme ruang
ukur dari $(\BbbR^2,\Sigma_2,\mu_2)$.

(e) Pemetaan $(x,y)\mapsto(x-y,y):\BbbR^2\to\BbbR^2$ merupakan automorfisme ruang
ukur dari $(\BbbR^2,\Sigma_2,\mu_2)$.

\cmmnt{\medskip

\noindent{\bf Catatan} Perlu saya catat bahwa (b), (d), dan (e) dapat
dipandang sebagai kasus-kasus khusus sederhana
dari Teorema 263A dalam bab berikutnya.   Namun demikian, menurut saya
tetap layak menuliskan bukti-bukti terpisah di sini, sebagian karena
kasus umum operator linear yang dibahas dalam 263A memerlukan beberapa
perangkat tambahan yang tidak diperlukan di sini, tetapi terlebih karena hasil di sini tidak ada
kaitannya dengan struktur {\it linear} dari $\Bbb R$ dan $\BbbR^2$;  hasil ini
semata-mata bergantung pada struktur {\it grup} dari $\Bbb R$, beserta
kaitan antara topologi dan ukurannya, dan argumen-argumen yang saya
berikan sekarang dapat disesuaikan dengan perumuman yang tepat ke grup
topologis abelian.

}%end of comment

\proof{{\bf (a)} Ini hanyalah keinvarian ukuran Lebesgue terhadap translasi,
yang dibahas dalam \S134.   Di sana saya menunjukkan bahwa jika $E\in\Sigma$
maka $E+a\in\Sigma$ dan $\mu(E+a)=\mu E$ (134Ab);  yakni, dengan menuliskan
$\phi(x)=x+a$, $\mu(\phi[E])$ ada dan sama dengan $\mu E$ untuk setiap
$E\in\Sigma$.   Tetapi tentu saja kita juga mempunyai

\Centerline{$\mu(\phi^{-1}[E])=\mu(E+(-a))=\mu E$}

\noindent untuk setiap $E\in\Sigma$, sehingga $\phi$ merupakan automorfisme.

\medskip

{\bf (b)} Pokoknya ialah bahwa $\mu^*(A)=\mu^*(-A)$ untuk setiap
$A\subseteq\Bbb R$.   \Prf\ (Saya mengikuti definisi-definisi dalam Jilid 1.)   Jika
$\epsilon>0$, terdapat suatu barisan $\sequencen{I_n}$ interval setengah terbuka
yang menutupi $A$ dengan $\sum_{n=0}^{\infty}\mu
I_n\le\mu^*A+\epsilon$.   Sekarang $-A\subseteq\bigcup_{n\in\Bbb N}(-I_n)$.
Tetapi jika $I_n=\coint{a_n,b_n}$ maka $-I_n=\ocint{-b_n,a_n}$, sehingga

\Centerline{$\mu^*(-A)\le\sum_{n=0}^{\infty}\mu(-I_n)
=\sum_{n=0}^{\infty}\max(0,-a_n-(-b_n))=\sum_{n=0}^{\infty}\mu I_n
\le\mu^*A+\epsilon$.}

\noindent Karena $\epsilon$ sebarang, $\mu^*(-A)\le\mu^*A$.   Tentu saja juga
$\mu^*A=\mu^*(-(-A))\le\mu^*(-A)$, sehingga $\mu^*(-A)=\mu^*A$.   \Qed

Ini berarti bahwa, dengan kali ini menetapkan $\phi(x)=-x$, $\phi$ merupakan
automorfisme dari struktur $(\Bbb R,\mu^*)$.   Tetapi karena $\mu$
didefinisikan dari $\mu^*$ melalui prosedur abstrak metode
\Caratheodory, $\phi$ juga harus merupakan automorfisme dari struktur $(\Bbb
R,\Sigma,\mu)$.

\medskip

{\bf (c)} Gabungkan (a) dan (b);  $x\mapsto a-x$ merupakan komposisi
automorfisme $x\mapsto -x$ dan $x\mapsto a+x$, dan
komposisi automorfisme sudah tentu merupakan automorfisme.

\medskip

{\bf (d)(i)} Tuliskan $\Tau$ untuk himpunan
$\{E:E\in\Sigma_2,\,\phi[E]\in\Sigma_2\}$, dengan kali ini
$\phi(x,y)=(x+y,y)$ untuk $x$, $y\in\Bbb R$, sehingga
$\phi:\BbbR^2\to\Bbb R^2$ merupakan suatu permutasi.
Maka $\Tau$ merupakan aljabar-$\sigma$, sebab merupakan
irisan aljabar-$\sigma$ $\Sigma_2$ dan
$\{E:\phi[E]\in\Sigma_2\}=\{\phi^{-1}[F]:F\in\Sigma_2\}$.   Selain itu,
$\mu_2E=\mu_2(\phi[E])$ untuk setiap $E\in\Tau$.   \Prf\  Menurut 252D,
kita mempunyai

\Centerline{$\mu_2E=\int\mu\{x:(x,y)\in E\}\mu(dy)$.}

\noindent Tetapi dengan menerapkan hasil yang sama pada $\phi[E]$ kita memperoleh

$$\eqalignno{\mu_2\phi[E]
&=\int\mu\{x:(x,y)\in\phi[E]\}\mu(dy)
=\int\mu\{x:(x-y,y)\in E\}\mu(dy)\cr
&=\int\mu(E^{-1}[\{y\}]+y)\mu(dy)
=\int\mu E^{-1}[\{y\}]\mu(dy)\cr
\noalign{\noindent (karena ukuran Lebesgue invarian terhadap translasi)}
&=\mu_2E. \text{  \Qed}\cr}$$

\medskip

\quad{\bf (ii)} Sekarang $\phi$ dan $\phi^{-1}$ jelas kontinu, sehingga
$\phi[G]$ terbuka, dan karena itu terukur, untuk setiap $G$ terbuka;
akibatnya semua himpunan terbuka harus termasuk dalam $\Tau$.   Karena $\Tau$ merupakan
aljabar-$\sigma$, ia memuat semua himpunan Borel.   Sekarang misalkan $E$ sebarang
himpunan terukur.   Maka terdapat himpunan Borel $H_1$, $H_2$ sedemikian sehingga
$H_1\subseteq E\subseteq H_2$ dan $\mu_2(H_2\setminus H_1)=0$ (134Fb).
Kita mempunyai $\phi[H_1]\subseteq\phi[E]\subseteq\phi[H_2]$ dan

\Centerline{$\mu_2(\phi[H_2]\setminus\phi[H_1])
=\mu_2\phi[H_2\setminus H_1]=\mu_2(H_2\setminus H_1)=0$.}

\noindent Jadi $\phi[E]\setminus\phi[H_1]$ harus terabaikan, sehingga
terukur, dan $\phi[E]=\phi[H_1]\cup(\phi[E]\setminus\phi[H_1])$
terukur.   Ini menunjukkan bahwa $\phi[E]$ terukur kapan pun $E$ terukur.

\medskip

\quad{\bf (iii)} Dengan mengulangi argumen yang sama menggunakan $-y$ sebagai pengganti
$y$, kita melihat bahwa $\phi^{-1}[E]$ terukur, dan
$\mu_2\phi^{-1}[E]=\mu_2E$, untuk setiap $E\in\Sigma_2$.
Jadi $\phi$ merupakan automorfisme dari struktur $(\BbbR^2,\Sigma_2,\mu_2)$.

\medskip

{\bf (e)} Tentu saja ini merupakan akibat langsung, baik dari bukti
(d) maupun dari pernyataan (d) itu sendiri, karena $(x,y)\mapsto(x-y,y)$ hanyalah
invers dari $(x,y)\mapsto(x+y,y)$.
}%end of proof of 255A

\leader{255B}{Akibat} (a) Jika $a\in\Bbb R$, maka untuk sebarang
fungsi bernilai kompleks $f$ yang didefinisikan pada suatu subhimpunan $\Bbb R$

\Centerline{$\int f(x)dx
=\int f(a+x)dx=\int f(-x)dx=\int f(a-x)dx$}

\noindent dalam arti bahwa jika salah satu integral ada, maka integral-integral
lainnya juga ada, dan semuanya sama.

(b) Jika $f$ merupakan fungsi bernilai kompleks yang didefinisikan pada suatu subhimpunan
$\BbbR^2$, maka

\Centerline{$\int f(x+y,y)d(x,y)=\int f(x-y,y)d(x,y)
=\int f(x,y)d(x,y)$}

\noindent dalam arti bahwa jika salah satu integral ada dan berhingga,
maka yang lain juga demikian, dan semuanya sama.

\cmmnt{
\leader{255C}{Catatan (a)}  Saya tidak yakin apakah seharusnya dianggap
`jelas' bahwa jika $(X,\Sigma,\mu)$, $(Y,\Tau,\nu)$ merupakan ruang ukur
dan $\phi:X\to Y$ suatu isomorfisme, maka untuk sebarang fungsi
$f$ yang didefinisikan pada suatu subhimpunan $Y$

\Centerline{$\int f(\phi(x))\mu(dx)=\int f(y)\nu(dy)$}

\noindent dalam arti bahwa jika yang satu terdefinisi maka yang lain pun terdefinisi, dan keduanya
sama.   Jika hal ini jelas, kejelasannya tentu bergantung
pada sifat definisi integrasi:
keterintegralan terhadap ukuran $\mu$ merupakan sesuatu yang
bergantung pada struktur $(X,\Sigma,\mu)$ dan tidak pada sifat lain apa pun dari
$X$.   Jika hal ini tidak jelas, maka ia merupakan deduksi mudah dari Teorema
235A di atas, yang diterapkan berturut-turut pada $\phi$ dan $\phi^{-1}$ serta pada bagian real
dan imajiner dari $f$.   Bagaimanapun juga, isomorfisme-isomorfisme dalam 255A
tepat merupakan isomorfisme yang diperlukan untuk membuktikan 255B.

\header{255Cb}{\bf (b)} Perhatikan bahwa dalam 255Bb saya menuliskan $\int f(x,y)d(x,y)$
untuk menekankan bahwa saya sedang meninjau integral $f$ terhadap
ukuran Lebesgue dua dimensi.   Fakta bahwa

\Centerline{$\int\bigl(\int f(x,y)dx\bigr)dy
=\int\bigl(\int f(x+y,y)dx\bigr)dy=\int\bigl(\int f(x-y,y)dx\bigr)dy$}

\noindent sebenarnya lebih mudah, karena merupakan akibat langsung dari
kesamaan $\int f(a+x)dx=\int f(x)dx$.   Tetapi penerapan hasil ini
sering kali pada dasarnya bergantung pada fakta bahwa fungsi-fungsi
$(x,y)\mapsto f(x+y,y)$, $(x,y)\mapsto f(x-y,y)$ terukur sebagai
fungsi dua variabel.

\header{255Cc}{\bf (c)} Saya langsung beralih ke fungsi bernilai
kompleks karena fungsi-fungsi ini diperlukan untuk penerapan dalam Bab
28.   Namun, jika
hal itu menimbulkan ketidaknyamanan bagi Anda, baik secara teknis maupun estetis, semua
gagasan teori ukuran dalam bagian ini sudah dapat ditemukan dalam
kasus real, dan mula-mula Anda mungkin ingin membacanya seolah-olah hanya bilangan real
yang terlibat.
}%end of comment

\leader{255D}{}\cmmnt{ Akibat lebih lanjut dari 255A akan berguna.

\medskip

\noindent}{\bf Akibat} Misalkan $f$ suatu fungsi bernilai kompleks yang didefinisikan
pada suatu subhimpunan $\Bbb R$.

(a) Jika $f$ terukur, maka fungsi-fungsi $(x,y)\mapsto f(x+y)$,
$(x,y)\mapsto f(x-y)$ terukur.

(b) Jika $f$ didefinisikan hampir di mana-mana dalam $\Bbb R$, maka fungsi-fungsi
$(x,y)\mapsto f(x+y)$, $(x,y)\mapsto f(x-y)$ didefinisikan hampir
di mana-mana dalam $\BbbR^2$.

\proof{ Dengan menuliskan $g_1(x,y)=f(x+y)$, $g_2(x,y)=f(x-y)$ kapan pun keduanya
terdefinisi, kita mempunyai

\Centerline{$g_1(x,y)=(f\otimes\chi\Bbb R)(\phi(x,y))$,
\quad $g_2(x,y)=(f\otimes\chi\Bbb R)(\phi^{-1}(x,y))$,}

\noindent dengan menuliskan $\phi(x,y)=(x+y,y)$ seperti dalam 255A(d-e), dan
$(f\otimes\chi\Bbb R)(x,y)=f(x)$, mengikuti notasi 253B.   Menurut
253C, $f\otimes\chi\Bbb R$ terukur jika $f$ terukur, dan didefinisikan
hampir di mana-mana jika $f$ demikian.   Karena $\phi$ merupakan automorfisme ruang
ukur, $(f\otimes\chi\Bbb R)\phi=g_1$ dan
$(f\otimes\chi\Bbb R)\phi^{-1}=g_2$ terukur, atau didefinisikan hampir
di mana-mana, jika $f$ demikian.
}%end of proof of 255D

\vleader{48pt}{255E}{Rumus dasar} Misalkan $f$ dan $g$ fungsi
bernilai kompleks yang terukur
dan didefinisikan hampir di mana-mana dalam $\Bbb R$.   Tuliskan $f*g$ untuk
fungsi yang didefinisikan oleh rumus

\Centerline{$(f*g)(x)=\int f(x-y)g(y)dy$}

\noindent kapan pun integral itu ada\cmmnt{ (terhadap
ukuran Lebesgue, tentu saja)} sebagai suatu bilangan kompleks.   Maka $f*g$ adalah
{\bf konvolusi} fungsi $f$ dan $g$.
\cmmnt{

Perhatikan bahwa} $\dom(|f|*|g|)=\dom(f*g)$, dan\cmmnt{ bahwa}
$|f*g|\le |f|*|g|$ di setiap titik domain bersama keduanya, untuk semua $f$ dan
$g$.

\cmmnt{\medskip

\noindent{\bf Catatan} Perhatikan bahwa di sini saya bersedia meninjau
konvolusi $f$ dan $g$ untuk sebarang anggota
$\eusm L^0_{\Bbb C}$, yakni ruang fungsi terukur bernilai kompleks yang
didefinisikan hampir di mana-mana, sekalipun
domain $f*g$ mungkin kosong.   }%end of comment

\leader{255F}{Sifat elementer (a)} Karena integrasi bersifat linear,
tentu kita mempunyai


\Centerline{$((f_1+f_2)*g)(x)=(f_1*g)(x)+(f_2*g)(x)$,}

\Centerline{$(f*(g_1+g_2))(x)=(f*g_1)(x)+(f*g_2)(x)$,}

\Centerline{$(cf*g)(x)=(f*cg)(x)=c(f*g)(x)$}

\noindent kapan pun ruas kanan rumus-rumus tersebut terdefinisi.

\spheader 255Fb Jika $f$ dan $g$ merupakan fungsi terukur bernilai kompleks
yang didefinisikan hampir di mana-mana dalam $\Bbb R$, maka
$f*g=g*f$\cmmnt{, dalam arti ketat bahwa keduanya mempunyai domain yang sama
dan nilai yang sama di setiap
titik domain bersama tersebut}.

\prooflet{\Prf\ Ambil $x\in\Bbb R$ dan terapkan 255Ba untuk melihat bahwa

$$\eqalign{(f*g)(x)
&=\int f(x-y)g(y)dy
=\int f(x-(x-y))g(x-y)dy\cr
&=\int f(y)g(x-y)dy
=(g*f)(x)\cr}$$

\noindent jika salah satunya terdefinisi.   \Qed}

\header{255Fc}{\bf (c)} Jika $f_1$, $f_2$, $g_1$, $g_2$ merupakan fungsi
terukur bernilai kompleks yang didefinisikan hampir di mana-mana dalam $\Bbb R$,
$f_1\eae f_2$ dan $g_1\eae g_2$, maka $f_1*g_1=f_2*g_2$.
\prooflet{\Prf\ Untuk setiap $x\in\Bbb R$ kita
akan mempunyai $f_1(x-y)=f_2(x-y)$ untuk hampir setiap $y\in\Bbb R$, menurut 255Ac.
Akibatnya $f_1(x-y)g_1(y)=f_2(x-y)g_2(y)$ untuk hampir setiap $y$,
dan $(f_1*g_1)(x)=(f_2*g_2)(x)$ dalam arti bahwa jika salah satunya
terdefinisi, maka yang lain juga terdefinisi, dan keduanya sama.\ \Qed}

Akibatnya, jika $u$, $v\in L^0_{\Bbb C}$, maka kita mempunyai suatu fungsi
$\theta(u,v)$
yang sama dengan $f*g$ kapan pun $f$, $g\in\eusm L^0_{\Bbb C}$
sedemikian sehingga $f^{\ssbullet}=u$ dan $g^{\ssbullet}=v$.
\cmmnt{Perhatikan bahwa $\theta(u,v)=\theta(v,u)$, dan
bahwa $\theta(u_1+u_2,v)$ memperluas $\theta(u_1,v)+\theta(u_2,v)$,
$\theta(cu,v)$ memperluas $c\theta(u,v)$ untuk semua $u$, $u_1$, $u_2$,
$v\in L^0_{\Bbb C}$ dan $c\in\Bbb C$.}

\leader{255G}{}\cmmnt{ Saya telah mengelompokkan 255Fa-255Fc karena
ketiganya hanya bergantung
pada gagasan sampai dengan 255Ac dan 255Ba.   Dengan menggunakan paruh kedua
255A dan 255B kita memperoleh hasil yang jauh lebih mendalam.   Saya mulai dengan apa yang tampaknya
merupakan hasil mendasar.

\medskip

\noindent}{\bf Teorema} Misalkan $f$, $g$, dan $h$ merupakan fungsi terukur
bernilai kompleks yang didefinisikan hampir di mana-mana dalam $\Bbb R$.

(a) Andaikan $\int h(x+y)f(x)g(y)d(x,y)$ ada dalam $\Bbb C$.
Maka

$$\eqalign{\int h(x)(f*g)(x)dx
&=\int h(x+y)f(x)g(y)d(x,y)\cr
&=\iint h(x+y)f(x)g(y)dxdy
=\iint h(x+y)f(x)g(y)dydx\cr}$$

\noindent asalkan
dalam ungkapan $h(x)(f*g)(x)$ kita menafsirkan produk itu sebagai $0$ jika
$h(x)=0$ dan $(f*g)(x)$ tidak terdefinisi.

(b) Jika, dengan penafsiran serupa terhadap $|h(x)|(|f|*|g|)(x)$, integral
$\int |h(x)|(|f|*|g|)(x)dx$ berhingga, maka
$\int h(x+y)f(x)g(y)d(x,y)$ ada dalam $\Bbb C$.

\proof{ Tinjau fungsi-fungsi

\Centerline{$k_1(x,y)=h(x)f(x-y)g(y)$,\quad $k_2(x,y)=h(x+y)f(x)g(y)$}

\noindent di mana pun fungsi-fungsi ini terdefinisi.   255D memberi tahu kita bahwa $k_1$ dan
$k_2$ terukur dan didefinisikan hampir di mana-mana.   Sekarang, dengan menetapkan
$\phi(x,y)=(x+y,y)$, kita mempunyai $k_2=k_1\phi$, sehingga

\Centerline{$\int k_1(x,y)d(x,y)=\int k_2(x,y)d(x,y)$}

\noindent jika salah satunya ada, menurut 255Bb.

Jika

\Centerline{$\int h(x+y)f(x)g(y)d(x,y)=\int k_2$}

\noindent ada, maka menurut teorema Fubini kita mempunyai

\Centerline{$\int k_2=\int k_1(x,y)d(x,y)
=\int(\int h(x)f(x-y)g(y)dy)dx$}

\noindent sehingga $\int h(x)f(x-y)g(y)dy$ ada hampir di mana-mana, yaitu,
$(f*g)(x)$ ada untuk hampir setiap $x$ sedemikian sehingga $h(x)\ne 0$;  dengan
penafsiran yang saya gunakan di sini, $h(x)(f*g)(x)$ ada hampir di mana-mana,
dan

$$\eqalign{\int h(x)(f*g)(x)dx
&=\int\bigl(\int h(x)f(x-y)g(y)dy\bigr)dx
=\int k_1\cr
&=\int k_2
=\int h(x+y)f(x)g(y)d(x,y)\cr
&=\iint h(x+y)f(x)g(y)dxdy
=\iint h(x+y)f(x)g(y)dydx\cr}$$

\noindent sekali lagi menurut teorema Fubini.

Jika (dengan penafsiran yang sama) $|h|\times(|f|*|g|)$ terintegralkan,

\Centerline{$|k_1(x,y)|=|h(x)||f(x-y)||g(y)|$}

\noindent terukur, dan

\Centerline{$\iint|h(x)||f(x-y)||g(y)|dydx
=\int|h(x)|(|f|*|g|)(x)dx$}

\noindent berhingga, sehingga menurut teorema Tonelli (252G, 252H), $k_1$ dan
$k_2$ terintegralkan.
}%end of proof of 255G

\leader{255H}{}\cmmnt{ Beberapa hasil standar kini mudah diperoleh.

\wheader{255H}{4}{2}{2}{24pt}

\noindent}{\bf Akibat} Jika $f$, $g$ merupakan fungsi bernilai kompleks
yang terintegralkan pada $\Bbb R$, maka $f*g$ terintegralkan,
dengan

\Centerline{$\int f*g =\int f\int g$,
\quad$\int|f*g|\le\int|f|\int|g|$.}

\proof{ Dalam 255G, tetapkan $h(x)=1$ untuk setiap $x\in\Bbb R$;  maka

\Centerline{$\int h(x+y)f(x)g(y)d(x,y)
=\int f(x)g(y)d(x,y)=\int f\int g$}

\noindent menurut 253D, sehingga

\Centerline{$\int f*g
=\int h(x)(f*g)(x)dx=\int h(x+y)f(x)g(y)d(x,y)=\int f\int g$,}

\noindent sebagaimana dinyatakan.   Sekarang

\Centerline{$\int|f*g|\le\int|f|*|g|=\int|f|\int|g|$.}
}%\end of proof of 255H

\leader{255I}{Akibat} Untuk sebarang fungsi terukur bernilai kompleks
$f$, $g$ yang didefinisikan hampir di mana-mana dalam $\Bbb R$, $f*g$ terukur dan
mempunyai domain terukur.

\proof{ Tetapkan $f_n(x)=f(x)$ jika $x\in\dom f$, $|x|\le n$ dan $|f(x)|\le n$,
dan $0$ di tempat lain dalam $\Bbb R$;  definisikan $g_n$ secara serupa dari $g$.   Maka
$f_n$ dan $g_n$ terintegralkan, $|f_n|\le|f|$ dan $|g_n|\le|g|$ hampir
di mana-mana, $f\eae\lim_{n\to\infty}f_n$ dan $g\eae\lim_{n\to\infty}g_n$.
Akibatnya, menurut Teorema Konvergensi Terdominasi Lebesgue,

$$\eqalign{(f*g)(x)
&=\int f(x-y)g(y)dy
=\int \lim_{n\to\infty}f_n(x-y)g_n(y)dy\cr
&=\lim_{n\to\infty}\int f_n(x-y)g_n(y)dy
=\lim_{n\to\infty}(f_n*g_n)(x)\cr}$$

\noindent untuk setiap $x\in\dom f*g$.   Namun $f_n*g_n$ terintegralkan,
sehingga terukur, untuk setiap $n$, jadi $f*g$ pasti terukur.

Adapun domain $f*g$,

$$\eqalign{x\in\dom(f*g)
&\iff\int f(x-y)g(y)dy\text{ terdefinisi dalam }\Bbb C\cr
&\iff\int|f(x-y)||g(y)|dy\text{ terdefinisi dalam }\Bbb R\cr
&\iff\int|f_n(x-y)||g_n(y)|dy\text{ terdefinisi dalam }\Bbb R
\text{ untuk setiap }n\cr
&\qquad\qquad\text{dan }
\sup_{n\in\Bbb N}\int|f_n(x-y)||g_n(y)|dy<\infty.\cr
}$$

\noindent Karena setiap $|f_n|*|g_n|$ terintegralkan, dan karenanya
terukur serta mempunyai domain terukur,

\Centerline{$\dom(f*g)=\{x:x\in\bigcap_{n\in\Bbb N}\dom(|f_n|*|g_n|),\,
\sup_{n\in\Bbb N}(|f_n|*|g_n|)(x)<\infty\}$}

\noindent terukur.
}%end of proof of 255I

\leader{255J}{Teorema} Misalkan $f$, $g$ dan $h$ fungsi terukur bernilai kompleks,
yang didefinisikan hampir di mana-mana dalam $\Bbb R$, sedemikian sehingga $f*g$ dan $g*h$
terdefinisi hampir di mana-mana.   Misalkan
$x\in\Bbb R$ sedemikian sehingga salah satu dari $(|f|*(|g|*|h|))(x)$,
$((|f|*|g|)*|h|)(x)$ terdefinisi dalam $\Bbb R$.   Maka $f*(g*h)$ dan
$(f*g)*h$ terdefinisi dan sama di $x$.

\proof{ Tetapkan $k(y)=f(x-y)$ ketika ruas ini terdefinisi, sehingga $k$
terukur dan didefinisikan hampir di mana-mana (255D).

\medskip

{\bf (a)} Jika $(|f|*(|g|*|h|))(x)$ terdefinisi, nilainya adalah
$\int|k(y)|(|g|*|h|)(y)dy$, sehingga menurut 255G kita mempunyai

\Centerline{$\int k(y)(g*h)(y)dy=\int k(y+z)g(y)h(z)d(y,z)$,}

\noindent yaitu,

$$\eqalign{(f*(g*h))(x)
&=\int f(x-y)(g*h)(y)dy
=\int k(y)(g*h)(y)dy\cr
&=\int k(y+z)g(y)h(z)d(y,z)
=\iint k(y+z)g(y)h(z)dydz\cr
&=\iint f(x-y-z)g(y)h(z)dydz
=\int (f*g)(x-z)h(z)dz\cr
&=((f*g)*h)(x).\cr}$$

\medskip

{\bf (b)} Jika $((|f|*|g|)*|h|)(x)$ terdefinisi, nilainya adalah

$$\eqalign{\int(|f|*|g|)(x-z)|h(z)|dz
&=\iint|f(x-z-y)||g(y)||h(z)|dydz\cr
&=\iint|k(y+z)||g(y)||h(z)|dydz.\cr}$$

\noindent Menurut 255D sekali lagi, $(y,z)\mapsto k(y+z)$ terukur, sehingga kita dapat
menerapkan teorema Tonelli untuk melihat bahwa $\int k(y+z)g(y)h(z)d(y,z)$ terdefinisi,
dan sama dengan $\int k(y)(g*h)(y)dy=(f*(g*h))(x)$ menurut 255Ga.   Di sisi
lain, menurut dua baris terakhir bukti (a),
$\int k(y+z)g(y)h(z)d(y,z)$ juga sama dengan $((f*g)*h)(x)$.
}%end of proof of  255J

\leader{255K}{}\cmmnt{ Saya rasa kita tidak akan memerlukan pembahasan menyeluruh
mengenai
persoalan kapan tepatnya $(f*g)(x)$ terdefinisi;  persoalan ini tampaknya
rumit.   Namun ada suatu kasus mendasar yang memungkinkan kita
memastikan bahwa $(f*g)(x)$ terdefinisi di mana-mana.

\medskip

\noindent}{\bf Proposisi} Misalkan $f$, $g$ fungsi terukur
bernilai kompleks
yang didefinisikan hampir di mana-mana dalam $\Bbb R$, dan
$f\in\eusm L^p_{\Bbb C}$,
$g\in\eusm L^q_{\Bbb C}$ dengan $p$, $q\in[1,\infty]$ dan
$\bover1p+\bover1q=1$\cmmnt{ (dengan menuliskan $\bover{1}{\infty}=0$ seperti
biasa)}.   Maka $f*g$ terdefinisi di mana-mana dalam $\Bbb R$, kontinu
seragam, dan

$$\eqalign{\sup_{x\in\Bbb R}|(f*g)(x)|
&\le\|f\|_p\|g\|_q\text{ jika }1<p<\infty,\,1<q<\infty,\cr
&\le\|f\|_1\esssup|g|\text{ jika }p=1,\,q=\infty,\cr
&\le\esssup|f|\cdot\|g\|_1\text{ jika }p=\infty,\,q=1.\cr}$$

\proof{{\bf (a)} (Untuk pengantar ruang $\eusm L^p$, lihat \S244.)
Untuk sebarang $x\in\Bbb R$, fungsi $f_x$, yang didefinisikan dengan menetapkan
$f_x(y)=f(x-y)$ kapan pun $x-y\in\dom f$, juga harus termasuk dalam
$\eusm L^p$, karena $f_x=f\phi$ untuk suatu automorfisme $\phi$ dari
ruang
ukur.   Akibatnya $(f*g)(x)=\int f_x\times g$ terdefinisi, dan modulusnya
paling besar $\|f\|_p\|g\|_q$ atau $\|f\|_1\esssup|g|$ atau
$\esssup|f|\cdot\|g\|_1$, menurut 244Eb/244Pb dan 243Fa/243K.

\medskip

{\bf (b)} Untuk melihat bahwa $f*g$ kontinu seragam, gunakan argumen berikut.
Misalkan mula-mula $p<\infty$.   Ambil $\epsilon>0$.   Ambil $\eta>0$
sedemikian sehingga $(2+2^{1/p})\|g\|_q\eta\le\epsilon$.
Maka ada fungsi kontinu terbatas $h:\Bbb R\to\Bbb C$ sedemikian
sehingga $\{x:h(x)\ne 0\}$ terbatas dan $\|f-h\|_p\le\eta$ (244Hb/244Pb);
ambil $M\ge 1$ sedemikian sehingga $h(x)=0$ kapan pun $|x|\ge M-1$.   Selanjutnya,
$h$ kontinu seragam, sehingga ada $\delta\in\ocint{0,1}$ sedemikian
sehingga $|h(x)-h(x')|\le M^{-1/p}\eta$ kapan pun $|x-x'|\le\delta$.

Misalkan $|x-x'|\le\delta$.   Dengan mendefinisikan $h_x(y)=h(x-y)$, seperti sebelumnya,
kita mempunyai

$$\eqalignno{\int|h_x-h_{x'}|^p
&=\int|h(x-y)-h(x'-y)|^pdy
=\int|h(t)-h(x'-x+t)|^pdt\cr
\displaycause{dengan menyubstitusikan $t=x-y$}
&=\int_{-M}^M|h(t)-h(x'-x+t)|^pdt\cr
\displaycause{karena $h(t)=h(x'-x+t)=0$ jika $|t|\ge M$}
&\le 2M(M^{-1/p}\eta)^p\cr
\displaycause{karena $|h(t)-h(x'-x+t)|\le M^{-1/p}\eta$ untuk setiap $t$}
&=2\eta^p.\cr}$$

\noindent Jadi $\|h_x-h_{x'}\|_p\le 2^{1/p}\eta$.   Di sisi lain,

\Centerline{$\int|h_{x}-f_{x}|^p
=\int|h(x-y)-f(x-y)|^pdy
=\int|h(y)-f(y)|^pdy$,}

\noindent sehingga $\|h_x-f_x\|_p=\|h-f\|_p\le\eta$, dan demikian pula
$\|h_{x'}-f_{x'}\|_p\le\eta$.   Jadi

\Centerline{$\|f_x-f_{x'}\|_p
\le\|f_x-h_x\|_p+\|h_x-h_{x'}\|_p+\|h_{x'}-f_{x'}\|_p
\le(2+2^{1/p})\eta$.}

\noindent Ini berarti bahwa

$$\eqalign{|(f*g)(x)-(f*g)(x')|
&=|\int f_x\times g-\int f_{x'}\times g|
=|\int(f_x-f_{x'})\times g|\cr
&\le\|f_x-f_{x'}\|_p\|g\|_q
\le(2+2^{1/p})\|g\|_q\eta
\le\epsilon.\cr}$$

\noindent Karena $\epsilon$ sebarang, $f*g$ kontinu seragam.

Argumen di sini mengandaikan bahwa $p$ berhingga.   Namun jika $p=\infty$ maka
$q=1$ berhingga, sehingga kita dapat menerapkan metode tersebut dengan $g$ menggantikan $f$ untuk
menunjukkan bahwa $g*f$ kontinu seragam, dan $f*g=g*f$ menurut 255Fb.
}%end of proof of 255K

\leader{255L}{Kasus berdimensi $r$} Saya telah menuliskan 255A-255K sebagai
teorema-teorema tentang ukuran Lebesgue pada $\Bbb R$.   Namun semuanya
berlaku\cmmnt{ sama baiknya} untuk ukuran Lebesgue pada $\BbbR^r$ bagi sebarang
$r\ge 1$\cmmnt{, dan perubahan
yang diperlukan begitu kecil sehingga saya rasa saya cukup
meminta Anda membaca kembali argumen-argumen tersebut, dengan mengubah setiap $\Bbb R$ menjadi
$\BbbR^r$, dan setiap $\BbbR^2$ menjadi $(\BbbR^r)^2$.   Dalam 255A
dan di tempat lain, ukuran $\mu_2$ harus dibaca sebagai ukuran Lebesgue
pada $\BbbR^{2r}$ atau sebagai ukuran produk pada $(\BbbR^r)^2$;
menurut 251N keduanya dapat diidentifikasi.   Ada perubahan sepele yang
diperlukan dalam bagian (b) bukti tersebut;  jika $I_n=\coint{a_n,b_n}$ maka

\Centerline{$\mu I_n=\mu(-I_n)
=\prod_{i=1}^r\max(0,\beta_{ni}-\alpha_{ni})$,}

\noindent dengan menuliskan $a_n=(\alpha_{n1},\ldots,\alpha_{nr}),\ b_n=(\beta_{n1},\ldots,\beta_{nr})$.   Dalam bukti
255I, fungsi-fungsi $f_n$ harus didefinisikan dengan menetapkan bahwa
$f_n(x)=f(x)$ jika $|f(x)|\le n$ dan $\|x\|\le n$, serta $0$ selainnya.

Karena itu, ketika mengutip hasil-hasil ini, saya tidak akan ragu untuk merujuk
paragraf-paragraf 255A-255K seolah-olah semuanya benar-benar telah dituliskan untuk
$r\ge 1$ secara umum}.

\leader{255M}{Kasus $\ocint{-\pi,\pi}$}\dvro{ {\bf (a)}}{ Gagasan yang sama
juga berlaku
untuk grup lingkaran $S^1$ dan interval $\ocint{-\pi,\pi}$, tetapi
di sini mungkin diperlukan penjelasan yang agak lebih panjang.

\spheader 255Ma Hal pertama yang harus ditetapkan ialah operasi
grup yang sesuai.}   Jika kita memandang $S^1$ sebagai himpunan
$\{z:z\in\Bbb C,\,|z|=1\}$, maka operasi grupnya adalah perkalian
kompleks, dan dalam
rumus-rumus di atas $x+y$ harus diganti dengan $xy$, sedangkan $x-y$ harus
diganti dengan $xy^{-1}$.   Pada interval $\ocint{-\pi,\pi}$, operasi grupnya
adalah $+_{2\pi}$, dengan, untuk $x$, $y\in\ocint{-\pi,\pi}$, saya menuliskan
$x+_{2\pi}y$ untuk salah satu dari $x+y$, $x+y+2\pi$, $x+y-2\pi$ yang termasuk dalam
$\ocint{-\pi,\pi}$.   \cmmnt{Untuk melihat bahwa ini memang suatu operasi
grup, salah satu
caranya ialah memperhatikan bahwa operasi ini berpadanan dengan perkalian pada $S^1$ jika kita
menggunakan bijeksi kanonik
$x\mapsto e^{ix}:\ocint{-\pi,\pi}\to S^1$;  cara lain ialah memperhatikan bahwa operasi ini
berpadanan dengan operasi pada grup hasil bagi $\Bbb R/2\pi\Bbb Z$.
Jadi, dalam penafsiran gagasan-gagasan 255A-255K ini, kita akan
mengganti $x+y$ dengan $x+_{2\pi}y$, $-x$ dengan $-_{2\pi}x$, dan $x-y$ dengan
$x-_{2\pi}y$, dengan

\Centerline{$-_{2\pi}x=-x$ jika $x\in\ooint{-\pi,\pi}$,
\quad$-_{2\pi}\pi=\pi$,}

\noindent dan $x-_{2\pi}y$ adalah salah satu dari $x-y$, $x-y+2\pi$, $x-y-2\pi$
yang termasuk dalam $\ocint{-\pi,\pi}$.
}%end of comment

\header{255Mb}{\bf (b)} Mengenai ukurannya, ukuran yang digunakan pada
$\ocint{-\pi,\pi}$ hanyalah ukuran Lebesgue.   \cmmnt{Perhatikan bahwa
karena
$\ocint{-\pi,\pi}$ terukur Lebesgue, tidak akan ada kerancuan
mengenai arti `subhimpunan terukur', karena subhimpunan yang terukur secara relatif
dari $\ocint{-\pi,\pi}$ sebenarnya diukur dengan
ukuran Lebesgue pada $\Bbb R$.   Kita juga dapat mengidentifikasi ukuran produk
pada $\ocint{-\pi,\pi}\times\ocint{-\pi,\pi}$ dengan ukuran subruang
yang diinduksi oleh ukuran Lebesgue pada $\BbbR^2$ (251R).}

Pada $S^1$, kita memerlukan ukuran padanan yang diinduksi oleh bijeksi
kanonik antara $S^1$ dan $\ocint{-\pi,\pi}$\cmmnt{, yang memang
sering disebut `ukuran Lebesgue pada $S^1$'}.   \cmmnt{(Kita akan melihat dalam
265E bahwa ukuran ini
juga sama dengan ukuran Hausdorff berdimensi satu pada $S^1$.)   Kita sudah sangat
dekat dengan taraf yang
memungkinkan kita secara wajar beralih ke $S^1$ dan ukuran ini (atau versi
ternormalisasinya, yang diperkecil dengan faktor $2\pi$, agar
$S^1$ menjadi ruang probabilitas).   Namun, teori elementer
deret Fourier, yang akan menjadi penerapan utama pembahasan ini dalam
jilid ini, umumnya dikerjakan pada interval dalam $\Bbb R$, sehingga
rumus-rumus yang didasarkan pada $\ocint{-\pi,\pi}$ lebih dekat dengan bentuk
standar.   Oleh karena itu, mulai sekarang saya akan menyajikan pembahasan
dalam istilah $\ocint{-\pi,\pi}$.}


\cmmnt{
\header{255Mc}{\bf (c)} Hasil yang berpadanan dengan 255A kini mempunyai
bentuk yang sedikit berbeda, jadi saya uraikan secara lengkap.
}

\leader{255N}{Teorema} Misalkan $\mu$ ukuran Lebesgue pada
$\ocint{-\pi,\pi}$ dan
$\mu_2$ ukuran Lebesgue pada $\ocint{-\pi,\pi}\times\ocint{-\pi,\pi}$;
tuliskan $\Sigma$, $\Sigma_2$ untuk domain-domainnya.

(a) Untuk sebarang $a\in\ocint{-\pi,\pi}$, peta $x\mapsto
a+_{2\pi}x:\ocint{-\pi,\pi}\to\ocint{-\pi,\pi}$ adalah suatu
automorfisme ruang ukur dari $(\ocint{-\pi,\pi},\Sigma,\mu)$.

(b) Peta $x\mapsto -_{2\pi}x:\ocint{-\pi,\pi}\to\ocint{-\pi,\pi}$ adalah
suatu automorfisme ruang ukur dari $(\ocint{-\pi,\pi},\Sigma,\mu)$.

(c) Untuk sebarang $a\in\ocint{-\pi,\pi}$, peta
$x\mapsto a-_{2\pi}x:\ocint{-\pi,\pi}\to\ocint{-\pi,\pi}$ adalah suatu
automorfisme ruang ukur dari $(\ocint{-\pi,\pi},\Sigma,\mu)$.

(d) Peta $(x,y)\mapsto(x+_{2\pi}y,y):\ocint{-\pi,\pi}^2\to
\ocint{-\pi,\pi}^2$ adalah suatu automorfisme ruang ukur dari
$(\ocint{-\pi,\pi}^2,
\ifdim\pagewidth>467pt\penalty-50\fi
\Sigma_2,
\ifdim\pagewidth>467pt\penalty-10\fi
\mu_2)$.

(e) Peta $(x,y)\mapsto(x-_{2\pi}y,y):\ocint{-\pi,\pi}^2\to
\ocint{-\pi,\pi}^2$ adalah suatu automorfisme ruang ukur dari
$(\ocint{-\pi,\pi}^2,
\ifdim\pagewidth>467pt\penalty-50\fi
\Sigma_2,
\ifdim\pagewidth>467pt\penalty-10\fi
\mu_2)$.

\proof{{\bf (a)} Tetapkan $\phi(x)=a+_{2\pi}x$.   Maka untuk sebarang
$E\subseteq\ocint{-\pi,\pi}$,

\Centerline{$\phi[E]
=((E+a)\cap\ocint{-\pi,\pi})
\cup(((E+a)\cap\ocint{\pi,3\pi})-2\pi)
\cup(((E+a)\cap\ocint{-3\pi,-\pi})+2\pi)$,}

\noindent dan ketiga himpunan ini saling lepas, sehingga

$$\eqalignno{\mu\phi[E]
&=\mu((E+a)\cap\ocint{-\pi,\pi})
+\mu(((E+a)\cap\ocint{\pi,3\pi})-2\pi)\cr
&\qquad\qquad\qquad\qquad\qquad\qquad
  +\mu(((E+a)\cap\ocint{-3\pi,-\pi})+2\pi)\cr
&=\mu_L((E+a)\cap\ocint{-\pi,\pi})
+\mu_L(((E+a)\cap\ocint{\pi,3\pi})-2\pi)\cr
&\qquad\qquad\qquad\qquad\qquad\qquad
  +\mu_L(((E+a)\cap\ocint{-3\pi,-\pi})+2\pi)\cr
\noalign{\noindent (dengan menuliskan $\mu_L$ untuk ukuran Lebesgue pada $\Bbb R$)}
&=\mu_L((E+a)\cap\ocint{-\pi,\pi})
+\mu_L((E+a)\cap\ocint{\pi,3\pi})
+\mu_L((E+a)\cap\ocint{-3\pi,-\pi})\cr
&=\mu_L(E+a)
=\mu_LE
=\mu E.\cr}$$

\noindent Demikian pula, $\mu\phi^{-1}[E]$ terdefinisi dan sama dengan $\mu E$
untuk setiap $E\in\Sigma$, sehingga $\phi$ merupakan automorfisme dari
$(\ocint{-\pi,\pi},
\ifdim\pagewidth>467pt\penalty-50\fi
\Sigma,
\ifdim\pagewidth>467pt\penalty-10\fi
\mu)$.

\medskip

{\bf (b)} Tentu saja bagian ini lebih singkat.   Dengan menetapkan $\phi(x)=-_{2\pi}x$ untuk
$x\in\ocint{-\pi,\pi}$,  kita mempunyai

$$\eqalign{\mu(\phi[E])
&=\mu(\phi[E]\cap\ooint{-\pi,\pi})
=\mu(-(E\cap\ooint{-\pi,\pi}))\cr
&=\mu_L(-(E\cap\ooint{-\pi,\pi}))
=\mu_L(E\cap\ooint{-\pi,\pi})\cr
&=\mu(E\cap\ooint{-\pi,\pi})
=\mu E\cr}$$

\noindent untuk setiap $E\in\Sigma$.

\medskip

{\bf (c)} Bagian ini cukup dengan menggabungkan (a) dan (b), seperti dalam
255A.

\medskip

{\bf (d)} Kita dapat berargumen seperti dalam (a), tetapi dengan sedikit rincian tambahan.
Jika $E\in\Sigma_2$, dan $\phi(x,y)=(x+_{2\pi}y,y)$ untuk $x$,
$y\in\ocint{-\pi,\pi}$, tetapkan $\psi(x,y)=(x+y,y)$ untuk $x$, $y\in\Bbb R$,
dan tuliskan $c=(2\pi,0)\in\BbbR^2$, $H=\ocint{-\pi,\pi}^2$, $H'=H+c$,
$H''=H-c$.   Maka untuk sebarang $E\in\Sigma_2$,

\Centerline{$\phi[E]=(\psi[E]\cap H)\cup((\psi[E]\cap H')-c)
\cup((\psi[E]\cap H'')+c)$,}

\noindent sehingga

$$\eqalignno{\mu_2\phi[E]
&=\mu_2(\psi[E]\cap H)
+\mu_2((\psi[E]\cap H')-c)
+\mu_2((\psi[E]\cap H'')+c)\cr
&=\mu_L(\psi[E]\cap H)
+\mu_L((\psi[E]\cap H')-c)
+\mu_L((\psi[E]\cap H'')+c)\cr
\noalign{\noindent (kali ini dengan menuliskan $\mu_L$ untuk ukuran Lebesgue pada
$\BbbR^2$)}
&=\mu_L(\psi[E]\cap H)
+\mu_L(\psi[E]\cap H')
+\mu_L(\psi[E]\cap H'')\cr
&=\mu_L\psi[E]
=\mu_LE
=\mu_2 E.\cr}$$

\noindent Dengan cara yang sama, $\mu_2(\phi^{-1}[E])=\mu_2E$ untuk setiap
$E\in\Sigma_2$, sehingga $\phi$ merupakan automorfisme dari
$(\ocint{-\pi,\pi}^2,\Sigma_2,\mu_2)$, sebagaimana diperlukan.

\medskip

{\bf (e)} Akhirnya, (e) hanyalah pernyataan ulang (d), seperti sebelumnya.
}%end of proof of 255N

\leader{255O}{Konvolusi pada $\ocint{-\pi,\pi}$}\cmmnt{ Dengan telah ditetapkannya
hasil
mendasar tersebut, argumen yang sama seperti dalam 255B-255K kini menghasilkan
hal berikut.}   Tuliskan $\mu$ untuk ukuran Lebesgue pada $\ocint{-\pi,\pi}$.


\header{255Oa}{\bf (a)} Misalkan $f$ dan $g$ fungsi terukur
bernilai kompleks
yang didefinisikan hampir di mana-mana dalam $\ocint{-\pi,\pi}$.   Tuliskan $f*g$
untuk fungsi yang didefinisikan oleh rumus

\Centerline{$(f*g)(x)=\int_{-\pi}^{\pi}f(x-_{2\pi}y)g(y)dy$}

\noindent kapan pun $x\in\ocint{-\pi,\pi}$ dan integral tersebut ada sebagai suatu
bilangan kompleks.   Maka $f*g$ adalah {\bf
konvolusi} dari fungsi-fungsi $f$ dan $g$.

\header{255Ob}{\bf (b)} Jika $f$ dan $g$ fungsi terukur bernilai kompleks
yang didefinisikan hampir di mana-mana dalam $\ocint{-\pi,\pi}$, maka
$f*g=g*f$.

\vspheader{60pt}255Oc Misalkan $f$, $g$ dan $h$ fungsi terukur bernilai kompleks
yang didefinisikan hampir di mana-mana dalam $\ocint{-\pi,\pi}$.   Maka

\quad{(i)}

\Centerline{$\int_{-\pi}^{\pi}h(x)(f*g)(x)dx
=\int_{\ocint{-\pi,\pi}^2}h(x+_{2\pi}y)f(x)g(y)d(x,y)$}

\noindent kapan pun ruas kanan ada dan berhingga, asalkan
dalam ungkapan $h(x)(f*g)(x)$ kita menafsirkan hasil kalinya sebagai $0$ jika
$h(x)=0$ dan $(f*g)(x)$ tidak terdefinisi.

\quad{(ii)} Jika, dengan penafsiran yang sama terhadap $|h(x)|(|f|*|g|)(x)$,
integral $\int_{-\pi}^{\pi}|h(x)|(|f|*|g|)(x)dx$
berhingga, maka $\int_{\ocint{-\pi,\pi}^2}h(x+_{2\pi}y)f(x)g(y)d(x,y)$ ada dalam
$\Bbb C$, sehingga sekali lagi kita mempunyai

\Centerline{$\int_{-\pi}^{\pi}h(x)(f*g)(x)dx
=\int_{\ocint{-\pi,\pi}^2}h(x+_{2\pi}y)f(x)g(y)d(x,y)$.}

\spheader 255Od Jika $f$, $g$ fungsi bernilai kompleks
yang terintegralkan pada $\ocint{-\pi,\pi}$, maka $f*g$
terintegralkan, dengan

\Centerline{$\int_{-\pi}^{\pi}f*g
=\int_{-\pi}^{\pi}f\int_{-\pi}^{\pi}g$,
\quad$\int_{-\pi}^{\pi}|f*g|
\le\int_{-\pi}^{\pi}|f|\int_{-\pi}^{\pi}|g|$.}

\spheader 255Oe Misalkan $f$, $g$, $h$ fungsi terukur bernilai kompleks
yang didefinisikan hampir di mana-mana dalam $\ocint{-\pi,\pi}$, sedemikian sehingga $f*g$
dan $g*h$ juga terdefinisi hampir di mana-mana.   Misalkan
$x\in\ocint{-\pi,\pi}$ sedemikian sehingga salah satu dari $(|f|*(|g|*|h|))(x)$,
$((|f|*|g|)*|h|)(x)$ terdefinisi dalam $\Bbb R$.   Maka $f*(g*h)$ dan
$(f*g)*h$ terdefinisi dan sama di $x$.

\spheader 255Of Misalkan $f\in\eusm L^p_{\Bbb C}(\mu)$,
$g\in\eusm L^{q}_{\Bbb C}(\mu)$ dengan $p$, $q\in[1,\infty]$ dan
$\bover1p+\bover1{q}=1$.   Maka $f*g$ terdefinisi
di mana-mana dalam $\ocint{-\pi,\pi}$, dan
$\sup_{x\in\ocint{-\pi,\pi}}|(f*g)(x)|\le\|f\|_p\|g\|_q$, dengan menafsirkan
$\|\,\,\|_{\infty}$ sebagai $\esssup |\,\,|$\cmmnt{, seperti dalam 255K}.

\exercises{
\leader{255X}{Latihan dasar $\pmb{>}$(a)}
%\sqheader 255Xa
Misalkan $f$, $g$ adalah fungsi bernilai kompleks yang terdefinisi hampir
di mana-mana dalam $\Bbb R$.   Tunjukkan bahwa untuk setiap
$x\in\Bbb R$, $(f*g)(x)=\int f(x+y)g(-y)dy$ jika salah satunya terdefinisi.
%255E

\sqheader 255Xb Misalkan $f$ dan $g$ adalah fungsi bernilai kompleks yang terdefinisi
hampir di mana-mana dalam $\Bbb R$.   (i) Tunjukkan bahwa jika $f$ dan $g$ adalah
fungsi genap, maka $f*g$ juga genap.   (ii) Tunjukkan bahwa jika $f$ genap dan $g$ ganjil
maka $f*g$ ganjil.   (iii) Tunjukkan bahwa jika $f$ dan $g$ ganjil maka $f*g$
genap.
%255E

\spheader 255Xc Misalkan $f$, $g$ adalah fungsi terukur bernilai real
yang terdefinisi hampir di mana-mana dalam $\BbbR^r$ dan sedemikian sehingga $f>0$
hampir di mana-mana, $g\ge 0$ hampir di mana-mana\ dan $\{x:g(x)>0\}$ tidak terabaikan.   Tunjukkan bahwa
$f*g>0$ di setiap titik dalam $\dom(f*g)$.
%255F

\sqheader 255Xd Misalkan $f:\Bbb R\to\Bbb C$ adalah fungsi terbatas yang
dapat didiferensialkan dan bahwa $f'$ terbatas.   Tunjukkan bahwa untuk setiap fungsi
bernilai kompleks yang terintegralkan
$g$ pada $\Bbb R$, $f*g$ dapat didiferensialkan dan
$(f*g)'=f'*g$ di setiap titik.   \Hint{123D.}
%255F

\spheader 255Xe Fungsi bernilai kompleks $g$ yang terdefinisi hampir di mana-mana
dalam $\Bbb R$ disebut {\bf terintegralkan secara lokal} jika $\int_a^bg$ terdefinisi dalam
$\Bbb C$ kapan pun $a<b$ dalam $\Bbb R$.   Misalkan $g$ adalah fungsi semacam itu
dan bahwa $f:\Bbb R\to\Bbb C$ adalah fungsi yang dapat didiferensialkan, dengan
turunan kontinu, sedemikian sehingga $\{x:f(x)\ne 0\}$ terbatas.   Tunjukkan
bahwa $(f*g)'=f'*g$ di setiap titik.
%255F 255Xd

\sqheader 255Xf Tetapkan $\phi_{\delta}(x)=\exp(-\bover1{\delta^2-x^2})$ jika
$|x|<\delta$, $0$ jika $|x|\ge\delta$, seperti dalam 242Xi.   Tetapkan
$\alpha_{\delta}=\int\phi_{\delta}$,
$\psi_{\delta}=\alpha_{\delta}^{-1}\phi_{\delta}$.   Misalkan $f$ adalah fungsi
bernilai kompleks yang terintegralkan secara lokal
pada $\Bbb R$.   (i) Tunjukkan bahwa $f*\psi_{\delta}$
adalah fungsi mulus yang terdefinisi di setiap titik dalam $\Bbb R$ untuk setiap
$\delta>0$.   (ii) Tunjukkan bahwa
$\lim_{\delta\downarrow 0}(f*\psi_{\delta})(x)=f(x)$ untuk hampir setiap
$x\in\Bbb R$.   \Hint{223Yg.}   (iii) Tunjukkan bahwa jika $f$ terintegralkan
maka $\lim_{\delta\downarrow 0}\int|f-f*\psi_{\delta}|=0$.   \Hint{gunakan
(ii) dan 245H(a-ii) {\it atau} perhatikan dahulu kasus $f=\chi[a,b]$ dan
gunakan 242O, dengan mencatat bahwa $\int|f*\psi_{\delta}|\le\int|f|$.}   (iv) Tunjukkan
bahwa jika $f$ kontinu seragam dan terdefinisi di setiap titik dalam $\Bbb R$
maka
$\lim_{\delta\downarrow 0}\sup_{x\in\Bbb R}|f(x)-(f*\psi_{\delta})(x)|=0$.
%255F 255Xe

\sqheader 255Xg Untuk $\alpha>0$, tetapkan
$g_{\alpha}(t)=\Bover1{\Gamma(\alpha)}t^{\alpha-1}$ untuk $t>0$, $0$ untuk
$t\le 0$.   Tunjukkan bahwa $g_{\alpha}*g_{\beta}=g_{\alpha+\beta}$ untuk semua
$\alpha$, $\beta>0$.   \Hint{252Yf.}
%255F

\sqheader 255Xh Misalkan $\mu$ adalah ukuran Lebesgue pada $\Bbb R$.
Untuk $u$, $v$,
$w\in L^0_{\Bbb C}=L^0_{\Bbb C}(\mu)$, katakan bahwa
$u*v=w$ jika $f*g$ terdefinisi hampir di mana-mana dan $(f*g)^{\ssbullet}=w$
kapan pun $f$, $g\in\eusm L^0_{\Bbb C}(\mu)$, $f^{\ssbullet}=u$ dan
$g^{\ssbullet}=v$.   (i) Tunjukkan bahwa $(u_1+u_2)*v=u_1*v+u_2*v$ kapan pun
$u_1$, $u_2$, $v\in L^0_{\Bbb C}$ dan $u_1*v$ serta $u_2*v$ terdefinisi dalam
pengertian ini.   (ii) Tunjukkan bahwa $u*v=v*u$ kapan pun $u$, $v\in L^0_{\Bbb C}$ dan
salah satu dari $u*v$ atau $v*u$ terdefinisi.   (iii) Tunjukkan bahwa jika $u$, $v$,
$w\in L^0_{\Bbb C}$, $u*v$ dan $v*w$ terdefinisi,
dan salah satu dari $|u|*(|v|*|w|)$ atau $(|u|*|v|)*|w|$
terdefinisi, maka $u*(v*w)=(u*v)*w$ terdefinisi dan sama.
%255I

\sqheader 255Xi Misalkan $\mu$ adalah ukuran Lebesgue pada $\Bbb R$.
(i) Tunjukkan bahwa $u*v$, sebagaimana didefinisikan dalam 255Xh, termasuk dalam $L^1_{\Bbb C}(\mu)$
kapan pun $u$, $v\in L^1_{\Bbb C}(\mu)$.
(ii) Tunjukkan bahwa $L^1_{\Bbb C}$ adalah
aljabar Banach komutatif terhadap $*$ (definisi: 2A4J).
%255Xh 255I

\spheader 255Xj(i) Tunjukkan bahwa jika $h$ adalah fungsi terintegralkan
pada $\BbbR^2$, maka   $(Th)(x)=\int h(x-y,y)dy$ ada untuk
hampir setiap $x\in\Bbb R$, dan bahwa $\int (Th)(x)dx=\int h(x,y)d(x,y)$.
(ii) Tuliskan $\mu_2$ untuk ukuran Lebesgue pada $\BbbR^2$, $\mu$ untuk
ukuran Lebesgue pada $\Bbb R$.   Tunjukkan bahwa ada operator linear
$\tilde T:L^1(\mu_2)\to L^1(\mu)$ yang didefinisikan dengan menetapkan
$\tilde T(h^{\ssbullet})=(Th)^{\ssbullet}$ untuk setiap fungsi terintegralkan
$h$ pada $\BbbR^2$.   (iii) Tunjukkan bahwa dalam bahasa 253E dan 255Xh,
$\tilde T(u\otimes v)=u*v$ untuk semua $u$, $v\in L^1(\mu)$.
%255Xh 255I 255Xi

\sqheader 255Xk  Untuk $\pmb{a}$, $\pmb{b}\in\Bbb C^{\Bbb Z}$ tetapkan
$(\pmb{a}*\pmb{b})(n)=\sum_{i\in\Bbb Z}\pmb{a}(n-i)\pmb{b}(i)$ kapan pun
$\sum_{i\in\Bbb Z}|\pmb{a}(n-i)\pmb{b}(i)|<\infty$.   Tunjukkan bahwa

\quad (i) $\pmb{a}*\pmb{b}=\pmb{b}*\pmb{a}$;

\quad (ii) $\sum_{i\in\Bbb Z}\pmb{c}(i)(\pmb{a}*\pmb{b})(i)
=\sum_{i,j\in\Bbb Z}\pmb{c}(i+j)\pmb{a}(i)\pmb{b}(j)$ jika
$\sum_{i,j\in\Bbb Z}|\pmb{c}(i+j)\pmb{a}(i)\pmb{b}(j)|<\infty$;

\quad (iii) jika $\pmb{a}$, $\pmb{b}\in\ell^1(\Bbb Z)$ maka
$\pmb{a}*\pmb{b}\in\ell^1(\Bbb Z)$ dan
$\|\pmb{a}*\pmb{b}\|_1\le\|\pmb{a}\|_1\|\pmb{b}\|_1$;

\quad (iv) Jika $\pmb{a}$, $\pmb{b}\in\ell^2(\Bbb Z)$ maka
$\pmb{a}*\pmb{b}\in\ell^{\infty}(\Bbb Z)$ dan
$\|\pmb{a}*\pmb{b}\|_{\infty}\le\|\pmb{a}\|_2\|\pmb{b}\|_2$;

\quad (v) jika $\pmb{a}$, $\pmb{b}$, $\pmb{c}\in\Bbb C^{\Bbb Z}$ dan
$(|\pmb{a}|*(|\pmb{b}|*|\pmb{c}|))(n)$ terdefinisi dengan baik, maka
$(\pmb{a}*(\pmb{b}*\pmb{c}))(n)=((\pmb{a}*\pmb{b})*\pmb{c})(n)$.
%255K

\leader{255Y}{Latihan lanjutan (a)}
%\spheader 255Ya
Misalkan $f$ adalah fungsi bernilai kompleks yang terintegralkan
pada $\Bbb R$.   (i) Misalkan $x$ adalah sebarang titik dalam himpunan Lebesgue dari $f$.
Tunjukkan bahwa untuk setiap $\epsilon>0$ terdapat $\delta>0$ sedemikian sehingga
$|f(x)-(f*g)(x)|\le\epsilon$ kapan pun $g:\Bbb R\to\coint{0,\infty}$ adalah
fungsi yang tak menurun pada $\ocint{-\infty,0}$, tak menaik
pada $\coint{0,\infty}$, dan memenuhi $\int g=1$ serta
$\int_{-\delta}^{\delta}g\ge 1-\delta$.   (ii) Tunjukkan bahwa untuk setiap
$\epsilon>0$ terdapat $\delta>0$ sedemikian sehingga $\|f-f*g\|_1\le\epsilon$
kapan pun $g:\Bbb R\to\coint{0,\infty}$ adalah fungsi yang
tak menurun pada $\ocint{-\infty,0}$, tak menaik pada
$\coint{0,\infty}$, dan memenuhi $\int g=1$ serta
$\int_{-\delta}^{\delta}g\ge 1-\delta$.
%255Xf 255F

\spheader 255Yb Misalkan $f$ adalah fungsi bernilai kompleks yang terintegralkan
pada $\Bbb R$.   Tunjukkan bahwa, untuk hampir setiap $x\in\Bbb R$,

\Centerline{$\lim_{a\to\infty}
  \Bover{a}{\pi}\int_{-\infty}^{\infty}
  \Bover{f(y)}{1+a^2(x-y)^2}dy$,
\quad$\lim_{a\to\infty}a\int_x^{\infty}f(y)e^{-a(y-x)}dy$,}

\Centerline{$\lim_{\sigma\downarrow 0}\Bover1{\sigma\sqrt{2\pi}}
  \int_{-\infty}^{\infty}f(y)e^{-(y-x)^2/2\sigma^2}dy$}

\noindent semuanya ada dan sama dengan $f(x)$.   \Hint{263G.}
%255Ya, 255Xf 255F

\spheader 255Yc
Tetapkan $f(x)=1$ untuk semua $x\in\Bbb R$, $g(x)=\Bover{x}{|x|}$
untuk $0<|x|\le 1$ dan $0$ selainnya, $h(x)=\tanh x$ untuk semua $x\in\Bbb R$.
Tunjukkan bahwa $f*(g*h)$ dan $(f*g)*h$ keduanya terdefinisi (dan konstan)
di setiap titik, dan keduanya berbeda.
%255J

\spheader 255Yd Bahas apa yang dapat terjadi jika, dalam konteks 255J, kita
mengetahui bahwa $(|f|*(|g|*|h|))(x)$ terdefinisi, tetapi tidak memiliki informasi tentang
domain $f*g$.
%255J

\spheader 255Ye Misalkan $p\in\coint{1,\infty}$
dan bahwa $f\in\eusm L^p_{\Bbb C}(\mu)$, dengan $\mu$ ukuran Lebesgue
pada $\BbbR^r$.   Untuk $a\in\BbbR^r$ tetapkan $(S_af)(x)=f(a+x)$ kapan pun
$a+x\in\dom f$.   Tunjukkan bahwa $S_af\in\eusm L^p_{\Bbb C}(\mu)$, dan bahwa
untuk setiap $\epsilon>0$ terdapat $\delta>0$ sedemikian sehingga
$\|S_af-f\|_p\le\epsilon$ kapan pun $|a|\le\delta$.
%255K

\spheader 255Yf Misalkan $p$, $q\in\ooint{1,\infty}$ dan
$\bover1p+\bover1q=1$.   Ambil $f\in\eusm L^p_{\Bbb C}(\mu)$ dan
$g\in\eusm L^q_{\Bbb C}(\mu)$, dengan $\mu$ ukuran Lebesgue pada
$\BbbR^r$.   Tunjukkan bahwa $\lim_{\|x\|\to\infty}(f*g)(x)=0$.
\Hint{gunakan 244Hb.}
%255K

\spheader 255Yg Ulangi 255Ye dan 255K, kali ini dengan mengambil $\mu$
sebagai ukuran Lebesgue pada $\ocint{-\pi,\pi}$, dan menetapkan
$(S_af)(x)=f(a+_{2\pi}x)$ untuk $a\in\ocint{-\pi,\pi}$;  tunjukkan bahwa dalam
versi baru 255K,
$(f*g)(\pi)=\lim_{x\downarrow -\pi}(f*g)(x)$.
%255K, 255Ye

\spheader 255Yh Misalkan $\mu$ adalah ukuran Lebesgue pada $\Bbb R$.   Untuk
$a\in\Bbb R$, $f\in\eusm L^0=\eusm L^0(\mu)$ tetapkan $(S_af)(x)=f(a+x)$
kapan pun $a+x\in\dom f$.

\quad (i) Tunjukkan bahwa $S_af\in\eusm L^0$ untuk setiap $f\in\eusm L^0$.

\quad  (ii) Tunjukkan bahwa kita mempunyai pemetaan $\tilde S_a:L^0\to L^0$ yang didefinisikan
dengan menetapkan $\tilde S_a(f^{\ssbullet})=(S_af)^{\ssbullet}$ untuk setiap
$f\in\eusm L^0$.

\quad  (iii) Tunjukkan bahwa $\tilde S_a$ adalah isomorfisme ruang Riesz dan
homeomorfisme untuk topologi konvergensi dalam ukuran;  selain itu,
bahwa $\tilde S_a(u\times v)=\tilde S_au\times\tilde S_av$ untuk semua $u$,
$v\in L^0$.

\quad  (iv) Tunjukkan bahwa $\tilde S_{a+b}=\tilde S_a\tilde S_b$ untuk semua $a$,
$b\in\Bbb R$.

\quad  (v) Tunjukkan bahwa $\lim_{a\to 0}\tilde S_au=u$ untuk topologi
konvergensi dalam ukuran, bagi setiap $u\in L^0$.

\quad  (vi) Tunjukkan bahwa jika $1\le p\le\infty$ maka $\tilde S_a\restr L^p$
adalah isomorfisme isometrik dari kisi Banach $L^p$.

\quad  (vii) Tunjukkan bahwa jika $p\in\coint{1,\infty}$ maka
$\lim_{a\to 0}\|\tilde S_au-u\|_p=0$ untuk setiap $u\in L^p$.

\quad  (viii) Tunjukkan bahwa jika $A\subseteq L^1$ terintegralkan secara seragam dan
$M\ge 0$, maka $\{\tilde S_au:u\in A,\,|a|\le M\}$ terintegralkan secara
seragam.

\quad (ix) Misalkan $u$, $v\in L^0$ sedemikian sehingga $u*v$ terdefinisi
dalam $L^0$ dalam pengertian 255Xh.   Tunjukkan bahwa
$\tilde S_a(u*v)=(\tilde S_au)*v=u*(\tilde S_av)$ untuk
setiap $a\in\Bbb R$.
%255Ye, 255K

\spheader 255Yi Buktikan 255Nd dari 255Na dengan metode yang digunakan untuk
membuktikan 255Ad dari 255Aa, alih-alih mengutip 255Ad.
%255N

\spheader 255Yj Misalkan $\mu$ adalah ukuran Lebesgue pada $\Bbb R$, dan
$\phi:\Bbb R\to\Bbb R$ fungsi konveks;  misalkan
$\bar\phi:L^0\to L^0=L^0(\mu)$ adalah operator terkait (lihat 241I).
Tunjukkan bahwa jika $u\in L^1=L^1(\mu)$, $v\in L^0$ sedemikian sehingga $u\ge 0$,
$\int u=1$ dan $u*v$, $u*\bar\phi(v)$ keduanya terdefinisi dalam pengertian
255Xh, maka $\bar\phi(u*v)\le u*\bar\phi(v)$.
\Hint{233I.}
%255Xh 255I

\spheader 255Yk Misalkan $\mu$ adalah ukuran Lebesgue pada $\Bbb R$, dan
$p\in[1,\infty]$.   Misalkan $f\in\eusm L^1_{\Bbb C}(\mu)$,
$g\in\eusm L^p_{\Bbb C}(\mu)$.   Tunjukkan bahwa
$f*g\in\eusm L^p_{\Bbb C}(\mu)$ dan bahwa
$\|f*g\|_p\le\|f\|_1\|g\|_p$.   \Hint{turunkan dari 255Yj, seperti dalam 244M.}
%255Yj 255K 255I

\spheader 255Yl Misalkan $p$, $q$, $r\in\ooint{1,\infty}$ dan bahwa
$\bover1p+\bover1q=1+\bover1r$.   Misalkan $\mu$ adalah ukuran Lebesgue pada
$\Bbb R$.   (i) Tunjukkan bahwa

\Centerline{$\int f\times g
\le\|f\|_p^{1-p/r}\|g\|_q^{1-q/r}(\int f^p\times g^q)^{1/r}$}

\noindent kapan pun $f$, $g\ge 0$ dan $f\in\eusm L^p(\mu)$,
$g\in\eusm L^q(\mu)$.   \Hint{tetapkan $p'=p/(p-1)$, dan seterusnya;  $f_1=f^{p/q'}$,
$g_1=g^{q/p'}$, $h=(f^p\times g^q)^{1/r}$.   Gunakan 244Xc untuk melihat bahwa
$\|f_1\times g_1\|_{r'}\le\|f_1\|_{q'}\|g_1\|_{p'}$, sehingga
$\int f_1\times g_1\times h\le\|f_1\|_{q'}\|g_1\|_{p'}\|h\|_r$.}   (ii)
Tunjukkan bahwa $f*g$ terdefinisi hampir di mana-mana\ dan bahwa
$\|f*g\|_r\le\|f\|_p\|g\|_q$ untuk semua $f\in\eusm L^p(\mu)$,
$g\in\eusm L^q(\mu)$.   \Hint{ambil $f$, $g\ge 0$.   Gunakan (i) untuk melihat bahwa
$(f*g)(x)^r\le\|f\|_p^{r-p}\|g\|_q^{r-q}\int f(y)^pg(x-y)^qdy$, sehingga
$\|f*g\|_r^r\le\|f\|_p^{r-p}\|g\|_q^{r-q}\int f(y)^p\|g\|_q^qdy$.}
(Ini adalah {\bf ketaksamaan Young}.)
%244Xc 255K

\spheader 255Ym Ulangi hasil-hasil bagian ini untuk grup
$(S^1)^r$, dengan $r\ge 2$, yang diberi ukuran produknya.
%+

\spheader 255Yn Misalkan $G$ adalah grup dan $\mu$ suatu ukuran
$\sigma$-hingga pada $G$ sedemikian sehingga ($\alpha$) untuk setiap $a\in G$,
pemetaan $x\mapsto ax$ adalah automorfisme dari $(G,\mu)$ ($\beta$) pemetaan
$(x,y)\mapsto (x,xy)$ adalah automorfisme dari $(G^2,\mu_2)$, dengan $\mu_2$
ukuran produk c.l.d.\ pada $G\times G$.   Untuk $f$,
$g\in\eusm L^0_{\Bbb C}(\mu)$ tuliskan $(f*g)(x)=\int f(y)g(y^{-1}x)dy$
kapan pun ini terdefinisi.   Tunjukkan bahwa
%+

\quad (i) jika $f$, $g$, $h\in\eusm L^0_{\Bbb C}(\mu)$ dan $\int
h(xy)f(x)g(y)d(x,y)$ terdefinisi dalam $\Bbb C$, maka $\int h(x)(f*g)(x)dx$
ada dan sama dengan $\int h(xy)f(x)g(y)d(x,y)$, asalkan
dalam ungkapan $h(x)(f*g)(x)$ kita menafsirkan hasil kali sebagai $0$ jika
$h(x)=0$ dan $(f*g)(x)$ tidak terdefinisi;

\quad (ii) jika $f$, $g\in\eusm L^1_{\Bbb C}(\mu)$ maka $f*g\in\eusm
L^1_{\Bbb C}(\mu)$ dan $\int f*g=\int f\int g$,
$\|f*g\|_1\le\|f\|_1\|g\|_1$;

\quad (iii) jika $f$, $g$, $h\in\eusm L^1_{\Bbb C}(\mu)$ maka
$f*(g*h)=(f*g)*h$.

\noindent (Lihat {\smc Halmos 50}, \S59.)

\spheader 255Yo Ulangi 255Yn untuk ukuran cacah pada sebarang grup $G$.
%+ 255Yn
}%end of exercises

\endnotes{
\Notesheader{255} Saya telah mencoba menyusun bagian ini sedemikian rupa
agar jelas bahwa dasar seluruh pembahasan di sini adalah 255A, dan
penerapan yang krusial adalah 255G.   Saya berharap bahwa jika dan ketika Anda mulai
mempelajari grup topologis umum (misalnya, dalam Bab 44), Anda
akan merasa mudah menelusuri
gagasan-gagasan tersebut dalam sebarang grup topologis abelian yang memungkinkan
Anda membuktikan suatu versi 255A.   Untuk grup nonabelian,
tentu saja, diperlukan kehati-hatian yang lebih besar, terutama karena dalam beberapa contoh penting
kita tidak lagi mempunyai $\mu\{x^{-1}:x\in E\}=\mu E$ untuk setiap $E$;
lihat 255Yn-255Yo untuk sedikit gambaran tentang apa yang dapat dilakukan tanpa menggunakan
gagasan topologis.

Pokok kritis dalam 255A adalah peralihan dari hasil berdimensi satu
dalam 255Aa-255Ac, yang tidak lain adalah invariansi terhadap translasi dan
refleksi dari
ukuran Lebesgue, ke hasil berdimensi dua dalam 255Ad-255Ae.   Dan
pusat argumen yang sesungguhnya, sebagaimana saya menyajikannya, adalah fakta bahwa
transformasi geser $\phi$ merupakan automorfisme dari struktur
$(\BbbR^2,\Sigma_2)$.   Perhitungan sebenarnya atas $\mu_2\phi[E]$,
dengan mengasumsikan bahwa himpunan itu terukur, adalah penerapan mudah dari teorema Fubini dan
Tonelli serta invariansi $\mu$ terhadap translasi.   Untuk
langkah inilah
kita mutlak memerlukan sifat-sifat topologis ukuran Lebesgue.
Barangkali saya perlu mengingatkan Anda bahwa fakta bahwa $\phi$ merupakan homeomorfisme
tidaklah cukup;  dalam 134I saya mendeskripsikan suatu homeomorfisme selang satuan
yang tidak melestarikan keterukuran, dan hal ini mudah diadaptasi
untuk menghasilkan homeomorfisme $\psi:\BbbR^2\to\BbbR^2$ sedemikian sehingga
$\psi[E]$ tidak selalu terukur untuk $E$ yang terukur.   Argumen
255A bergantung pada hubungan khusus di antara ketiga unsur berikut:
ukuran, topologi, dan struktur grup dari $\Bbb R$.

Saya telah mengemukakan beberapa komentar mengenai apa yang semestinya, atau tidak semestinya,
dianggap `jelas' (255C).   Namun barangkali dapat saya tambahkan bahwa hasil seperti
255B dan pernyataan kemudian, dalam pembuktian 255K, bahwa versi tercermin
dari suatu fungsi dalam $\eusm L^p$ juga termasuk dalam $\eusm L^p$, hanya dapat menjadi konsekuensi
sepele dari hasil seperti 255A jika setiap langkah dalam konstruksi
integral dilakukan dalam konteks abstrak ruang ukur umum.
Meskipun di sini kita bekerja secara eksklusif dengan integral Lebesgue,
argumen akan menjadi tidak dapat dipercaya jika pada tahap mana pun dalam
definisi integral kita bahkan menyebutkan bahwa yang sedang kita pikirkan adalah
ukuran Lebesgue.   Saya mengajukan hal ini sebagai alasan kuat untuk mendefinisikan
`integrasi' pada ruang ukur abstrak sejak awal, sebagaimana saya lakukan dalam
Jilid 1.   Bahkan, saya menyarankan bahwa secara umum dalam matematika murni ada
alasan baik untuk menuangkan argumen ke dalam bentuk-bentuk yang sesuai dengan
argumen itu sendiri.

Saya menulis buku ini bagi pembaca yang berminat pada pembuktian, dan seperti
di bagian lain saya telah menuliskan pembuktian-pembuktian bagian ini secara terperinci.   Namun
kebanyakan dari kita merasa berguna untuk menelaah sejumlah
materi dalam modus `kalkulus lanjut', yang saya maksudkan sebagai memulai dengan
rumus seperti

\Centerline{$(f*g)(x)=\int f(x-y)g(y)dy$,}

\noindent lalu menurunkan konsekuensi melalui manipulasi formal,
misalnya

\Centerline{$\int h(x)(f*g)(x)dx=\iint h(x)f(x-y)g(y)dydx
=\iint h(x+y)f(x)g(y)dydx$,}

\noindent tanpa terlebih dahulu mempersoalkan keberlakuan tepat dari
rumus-rumus tersebut.   Dalam beberapa hal, pendekatan berbasis rumus ini dapat
lebih setia pada struktur pokok bahasan daripada analisis cermat yang
biasa saya sajikan.   Hipotesis persis yang diperlukan agar teorema-teorema
benar secara ketat tentunya bersifat sekunder, dalam konteks seperti
bagian ini, terhadap pola yang dibentuk oleh keseluruhan teorema,
yang dapat diungkapkan secara memadai dan elegan dalam rumus-rumus
lugas.   Tentu saja saya tetap berkeras bahwa kita tidak dapat benar-benar
menghargai strukturnya, ataupun menggunakannya dengan aman, tanpa menguasai gagasan
dalam pembuktian -- dan sebagaimana telah saya katakan di tempat lain, saya percaya bahwa penguasaan
gagasan selalu mencakup penguasaan perincian formal, sekurang-kurangnya dalam arti
mampu merekonstruksinya dengan cukup lancar
saat diminta.

Sepanjang uraian utama bagian ini, saya telah bekerja dengan
fungsi alih-alih kelas ekuivalensi fungsi.    Namun semua
hasil di sini mempunyai penafsiran yang sangat penting bagi teori
`ruang fungsi' dalam Bab 24.   Dalam 255Xh dan latihan-latihan
sesudahnya, saya telah menunjukkan definisi konvolusi sebagai operator
dari suatu subhimpunan $L^0\times L^0$ ke $L^0$.
Satu hal yang menarik
adalah bahwa jika $u$, $v\in L^0$ maka $u*v$ dapat ditafsirkan sebagai
suatu {\it fungsi}, bukan sebagai anggota $L^0$ (255Fc).
Dengan demikian 255H dapat dipandang sebagai menyatakan bahwa $u*v\in\eusm L^1$
untuk $u$, $v\in L^1$.   Kita tidak dapat persis mengatakan bahwa konvolusi adalah
operator bilinear dari $L^1\times L^1$ ke $\eusm L^1$, karena $\eusm L^1$,
sebagaimana saya definisikan,
secara ketat bukanlah ruang linear.   Jika kita menginginkan operator bilinear,
maka kita harus memandang konvolusi sebagai fungsi dari
$L^1\times L^1$ ke $L^1$.   Namun ketika kita memandang konvolusi sebagai
fungsi pada $L^2\times L^2$, misalnya, maka fungsi-fungsi $u*v$ kita
terdefinisi di setiap titik (255K), dan bahkan merupakan fungsi kontinu yang lenyap
di $\infty$ (255Ye-255Yf).   Jadi dalam kasus ini tampaknya lebih tepat
memandang konvolusi sebagai operator bilinear dari $L^2\times L^2$ ke
suatu ruang fungsi kontinu, dan bukan sebagai operator dari
$L^2\times L^2$ ke $L^{\infty}$.   Untuk contoh konvolusi menarik
yang tidak secara alami dapat direpresentasikan dalam bentuk operator
pada ruang $L^p$, lihat 255Xg.

Karena konvolusi bertindak sebagai operator bilinear kontinu dari
$L^1(\mu)\times L^1(\mu)$ ke $L^1(\mu)$, dengan $\mu$ ukuran Lebesgue
pada $\Bbb R$, Teorema 253F memberi tahu kita bahwa konvolusi harus berkorespondensi dengan operator linear
dari $L^1(\mu_2)$ ke $L^1(\mu)$, dengan $\mu_2$ ukuran Lebesgue
pada $\BbbR^2$.  Inilah operator $\tilde T$ dari 255Xj.

Sejauh ini dalam catatan ini saya telah menulis seolah-olah kita hanya membahas
ukuran Lebesgue pada $\Bbb R$.   Namun banyak penerapan gagasan-gagasan tersebut
melibatkan $\BbbR^r$ atau $\ocint{-\pi,\pi}$ atau $S^1$.   Peralihan ke
$\BbbR^r$ semestinya elementer.   Peralihan ke $S^1$ memang memerlukan
perumusan ulang hasil dasar 255A/255N.   Juga semestinya jelas
bahwa tidak akan ada kesulitan baru dalam beralih ke $\ocint{-\pi,\pi}^r$
atau $(S^1)^r$.   Selain itu, kita juga dapat menjalani seluruh teori untuk
grup $\Bbb Z$ dan $\Bbb Z^r$, dengan ukuran yang sesuai kini adalah
ukuran cacah, sehingga $L^0_{\Bbb C}$ menjadi teridentifikasi dengan
$\Bbb C^{\Bbb Z}$ atau $\Bbb C^{\Bbb Z^r}$ (255Xk, 255Yo).


}%end of notes

\discrpage
